% !TEX root =./ECE444_Practice_main.tex
%
% L12 practice set (Pattern Measurement Theory) -- range length from the pi/8
% criterion, quiet-zone and absorber reasoning, comparison-method gain, and the
% near-field scan and pattern-cut bookkeeping behind the midterm project.

\fullwidth{\textbf{Learning Objective 2.5: } I can explain the theory behind antenna
pattern measurement, including anechoic chambers, near-field to far-field
transformations, and standard gain horns.
\\[4pt]\normalfont\small Show your work. A \textbf{Documentation} line at the end
is required (write \textbf{None} if you did not collaborate).}

\begin{questions}

\question \textbf{Sizing the range.} A $0.75\ \m$ aperture antenna is to be
pattern-tested at $8\ \ghz$. Take $c = 3\ee{8}\ \m/\text{s}$ throughout.
\begin{parts}

	\begin{minipage}[t]{0.58\linewidth}
		\part Find the wavelength and the minimum far-field range length for
		this antenna.
		\begin{solution}[1.2in]
		\eq{ \lambda &= \frac{c}{f} = \frac{3\ee{8}}{8\ee{9}} = 0.0375\ \m \\
		     r_{ff} &= \frac{2D^2}{\lambda} = \frac{2\lp 0.75 \rp^2}{0.0375}
		            = \frac{1.125}{0.0375} = 30.0\ \m }
		\end{solution}
	\end{minipage}\hfill%
	\begin{minipage}[t]{0.37\linewidth}
		\raggedleft \vspace{12pt}
		$r_{ff}$ = \ansbox[1.3in]{$30.0\ \m$}
	\end{minipage}

	\begin{minipage}[t]{0.58\linewidth}
		\part The only chamber available puts the source $12\ \m$ from the
		positioner. Find the edge phase error, in degrees, that this range
		imposes across the antenna.
		\begin{solution}[1.3in]
		\eq{ \Delta\phi_{max} &= \frac{\pi D^2}{4 \lambda r}
		     = \frac{\pi \lp 0.75 \rp^2}{4 \lp 0.0375 \rp \lp 12 \rp}
		     = \frac{1.767}{1.80} = 0.982\ \text{rad} \\
		     &= 56.3\mydeg }
		Equivalently, $r/r_{ff} = 12/30 = 0.40$, and the edge error scales as the
		inverse of that ratio: $22.5\mydeg / 0.40 = 56.3\mydeg$.
		\end{solution}
	\end{minipage}\hfill%
	\begin{minipage}[t]{0.37\linewidth}
		\raggedleft \vspace{12pt}
		$\Delta\phi_{max}$ = \ansbox[1.3in]{$56.3\mydeg$}
	\end{minipage}

	\begin{minipage}[t]{0.58\linewidth}
		\part A colleague must test a $2.4\ \m$ reflector at $35\ \ghz$. Find
		the far-field range length her measurement would require.
		\begin{solution}[1.2in]
		\eq{ \lambda &= \frac{3\ee{8}}{35\ee{9}} = 8.57\ee{-3}\ \m \\
		     r_{ff} &= \frac{2 \lp 2.4 \rp^2}{8.57\ee{-3}} = \frac{11.52}{8.57\ee{-3}} \\
		            &= 1344\ \m \approx 1.34\ \km }
		\end{solution}
	\end{minipage}\hfill%
	\begin{minipage}[t]{0.37\linewidth}
		\raggedleft \vspace{12pt}
		$r_{ff}$ = \ansbox[1.3in]{$1.34\ \km$}
	\end{minipage}

	\part Your colleague has no 1.3 km of chamber and no budget for a
	1.3 km outdoor range. Name the two measurement approaches that solve her
	problem, and state in one sentence what each one trades away to do it.
		\begin{solution}[1.4in]
		A \textbf{compact range} puts a precision offset-parabolic reflector in
		the room and collimates the feed's spherical wave into a plane wave over
		a short distance; it trades away money and simplicity, since the usable
		quiet zone is only about 50--60\% of the reflector aperture (so the
		reflector must be far larger than the antenna) and the rim must be
		serrated or rolled to keep edge diffraction out of the test volume.

		A \textbf{near-field scanner} measures amplitude and phase on a surface a
		few wavelengths from the aperture and transforms the result to the far
		field; it trades away speed and simplicity, since it needs a
		phase-coherent receiver, half-wavelength sampling over a large surface,
		and probe compensation.
		\end{solution}

\end{parts}

\newpage

\question \textbf{Reading a suspect pattern.} A short range does not fail
loudly. It returns a pattern that still looks like a pattern.
\begin{parts}

	\part In one or two sentences, explain physically why a source at finite
	range produces a phase error across the antenna under test that grows as the
	\emph{square} of the distance from the aperture centre.
		\begin{solution}[1.2in]
		The incident wavefront is a sphere of radius $r$ centred on the source,
		not a plane. A point a distance $x$ off the aperture centre is
		$\sqrt{r^2 + x^2} \approx r + x^2/2r$ away from the source, so the extra
		path -- and therefore the extra phase lag -- is $x^2/2r$. The linear term
		vanishes because the geometry is symmetric about boresight, leaving the
		quadratic term as the leading error.
		\end{solution}

	\begin{minipage}[t]{0.58\linewidth}
		\part An antenna with $D/\lambda = 25$ is measured at half its far-field
		distance. Find the edge phase error in degrees.
		\begin{solution}[1.2in]
		The far-field criterion is exactly $\pi/8 = 22.5\mydeg$ of edge error,
		and the error scales as $1/r$:
		\eq{ \Delta\phi_{max} = \frac{22.5\mydeg}{r/r_{ff}}
		     = \frac{22.5\mydeg}{0.5} = 45\mydeg }
		Note that this does not depend on $D/\lambda$ at all -- only on how short
		the range is relative to $2D^2/\lambda$.
		\end{solution}
	\end{minipage}\hfill%
	\begin{minipage}[t]{0.37\linewidth}
		\raggedleft \vspace{12pt}
		$\Delta\phi_{max}$ = \ansbox[1.3in]{$45\mydeg$}
	\end{minipage}

	\part In that same measurement the first null reads $-16\db$ instead of
	going deep, but the measured half-power beamwidth agrees with prediction to
	better than 2\%. Explain, in one or two sentences, how both statements can be
	true at once.
		\begin{solution}[1.2in]
		Quadratic phase error is a defocus, and a defocus destroys the precise
		cancellation that creates a null long before it changes the width of the
		main beam. The nulls depend on complete destructive interference across
		the whole aperture, so a $45\mydeg$ edge error fills them to about
		$-16\db$; the half-power points sit where the pattern is still dominated
		by the aperture's size, which the phase error has not changed. This is
		exactly why $2D^2/\lambda$ is a main-beam and gain criterion, not a
		sidelobe criterion.
		\end{solution}

	\part A customer needs the $-25\db$ sidelobe of this antenna verified. Using
	the fact that a range at $2D^2/\lambda$ fills the first null to about
	$-22\db$, at $5D^2/\lambda$ to about $-30\db$, and at $10D^2/\lambda$ to about
	$-36\db$, recommend a range length and defend it in one sentence.
		\begin{solution}[1.3in]
		Recommend $r \geq 10D^2/\lambda$, i.e.\ five times the usual far-field
		distance. At $2D^2/\lambda$ the range's own artifact sits at $-22\db$,
		which is \emph{above} the $-25\db$ feature being measured -- the
		measurement would be reporting the range, not the antenna. At
		$5D^2/\lambda$ the artifact is $-30\db$, only $5\db$ below the sidelobe,
		which is still worth several dB of error. At $10D^2/\lambda$ the artifact
		is $-36\db$, a comfortable $11\db$ below the feature, and the sidelobe
		level can be quoted with confidence.
		\end{solution}

\end{parts}

\newpage

\question \textbf{Specifying the chamber.} An anechoic chamber is not a room
with no reflections. It is a room with reflections below a stated level.
\begin{parts}

	\part Absorber reflectivity is quoted at normal incidence and degrades at
	grazing incidence. Explain why, and state which chamber surfaces this makes
	the critical ones.
		\begin{solution}[1.2in]
		The pyramids work as a gradual impedance taper: at normal incidence a
		wave travels a long way down into the lossy foam before it can reflect,
		so it is absorbed. At grazing incidence the wave sees the pyramid tips as
		a nearly continuous surface and travels only a short distance in the
		lossy material, so it reflects far more strongly. The critical surfaces
		are therefore the \textbf{side walls, ceiling, and floor} in the specular
		region between the source and the antenna under test, where the stray
		rays strike at a slant.
		\end{solution}

	\begin{minipage}[t]{0.58\linewidth}
		\part A single stray reflection arrives $40\db$ below the main beam.
		Find the worst-case error, in dB, that it adds to a measurement of
		(i) the main beam peak and (ii) a $-30\db$ sidelobe.
		\begin{solution}[1.4in]
		The stray adds in and out of phase with the wanted signal. Work with
		voltage ratios. For the main beam, the ratio is
		$10^{-40/20} = 0.0100$:
		\eq{ 20\ \mylog{1 \pm 0.0100} = +0.09\db \text{ / } -0.09\db }
		For a $-30\db$ sidelobe the stray is only $10\db$ below the wanted
		signal, so the ratio is $10^{-10/20} = 0.316$:
		\eq{ 20\ \mylog{1 \pm 0.316} = +2.4\db \text{ / } -3.3\db }
		\end{solution}
	\end{minipage}\hfill%
	\begin{minipage}[t]{0.37\linewidth}
		\raggedleft \vspace{12pt}
		main beam = \ansbox[1.1in]{$\pm 0.09\db$}
		\\[10pt]
		sidelobe = \ansbox[1.1in]{$+2.4/-3.3\db$}
	\end{minipage}

	\begin{minipage}[t]{0.58\linewidth}
		\part You must verify a $-35\db$ sidelobe to within $\pm 1\db$. Find the
		chamber reflectivity, relative to the main beam, that this requires.
		\begin{solution}[1.4in]
		Let $x$ be the stray-to-signal voltage ratio at the sidelobe. The
		negative-going error is the binding one:
		\eq{ 20\ \mylog{1 - x} &\geq -1\db \\
		     1 - x &\geq 10^{-1/20} = 0.891 \\
		     x &\leq 0.109 \quad \lp -19.3\db \rp }
		The stray must therefore sit $19.3\db$ below the sidelobe, which is
		$19.3\db$ below $-35\db$:
		\eq{ -35\db - 19.3\db = -54.3\db }
		Specify $-55\db$ over the band and quiet zone.
		\end{solution}
	\end{minipage}\hfill%
	\begin{minipage}[t]{0.37\linewidth}
		\raggedleft \vspace{12pt}
		reflectivity = \ansbox[1.3in]{$\approx -55\db$}
	\end{minipage}

	\part Your antenna under test is $1.5\ \m$ across. The chamber datasheet
	advertises a $1.2\ \m$ quiet zone at $-45\db$. State the problem in one
	sentence, and say what happens to the measurement if you proceed anyway.
		\begin{solution}[1.1in]
		The antenna does not fit inside the volume the chamber is certified for,
		so its outer $0.15\ \m$ on each side is illuminated by fields that carry
		no reflectivity guarantee at all. Proceeding means the aperture edges --
		the region that sets the sidelobe structure -- are being measured in an
		ordinary room, and any sidelobe or null result from the run is
		indefensible even though the main beam and gain will still look fine.
		\end{solution}

\end{parts}

\newpage

\question \textbf{Two ways to an absolute number.} Pattern shape is free; gain
costs you a calibration.
\begin{parts}

	\begin{minipage}[t]{0.58\linewidth}
		\part At $12\ \ghz$ a standard gain horn with a calibrated gain of
		$20.4\dbi$ receives $-41.2\dbm$. Swapping in the antenna under test,
		with nothing else touched, gives $-29.8\dbm$. Find the gain of the
		antenna under test.
		\begin{solution}[1.2in]
		\eq{ \Delta P &= P_{AUT} - P_{SGH} = -29.8 - \lp -41.2 \rp \\
		             &= +11.4\db \\
		     G_{AUT} &= G_{SGH} + \Delta P = 20.4 + 11.4 = 31.8\dbi }
		\end{solution}
	\end{minipage}\hfill%
	\begin{minipage}[t]{0.37\linewidth}
		\raggedleft \vspace{12pt}
		$G_{AUT}$ = \ansbox[1.3in]{$31.8\dbi$}
	\end{minipage}

	\part The transmit power, the source antenna's gain, the range length, and
	the cable losses were all unknown in part (a), yet the answer is exact.
	Explain in one or two sentences why they do not matter.
		\begin{solution}[1.1in]
		Write Friis for each of the two measurements. Every one of those
		quantities appears identically in both, because nothing but the antenna
		on the positioner changed, so all of them cancel in the difference
		$P_{AUT} - P_{SGH}$. The comparison method measures a \emph{ratio} of
		gains, and the horn's calibration supplies the only absolute number in
		the room.
		\end{solution}

	\begin{minipage}[t]{0.58\linewidth}
		\part No standard is available, so you set up two identical antennas
		facing each other at $R = 15\ \m$ and $10\ \ghz$. With $P_t = +10\dbm$
		you measure $P_r = -30.0\dbm$. Find the gain of one antenna.
		\begin{solution}[1.5in]
		Friis in dB, with $G_t = G_r = G$:
		\eq{ \lp P_r - P_t \rp_{dB} = 2G + 20\ \mylog{\frac{\lambda}{4\pi R}} }
		With $\lambda = 0.03\ \m$:
		\eq{ \frac{\lambda}{4\pi R} &= \frac{0.03}{4\pi \lp 15 \rp} = 1.592\ee{-4} \\
		     20\ \mylog{1.592\ee{-4}} &= -76.0\db }
		\eq{ 2G &= \lp -30.0 - 10.0 \rp - \lp -76.0 \rp = 36.0\db \\
		     G &= 18.0\dbi }
		\end{solution}
	\end{minipage}\hfill%
	\begin{minipage}[t]{0.37\linewidth}
		\raggedleft \vspace{12pt}
		$G$ = \ansbox[1.3in]{$18.0\dbi$}
	\end{minipage}

	\part Back in part (a): the standard gain horn is well matched, but the
	antenna under test has a VSWR of 2.0. Find the size of the resulting error,
	say which direction it pushes the answer, and state how you would report the
	result.
		\begin{solution}[1.5in]
		VSWR $= 2.0$ gives $\vert \Gamma \vert = \lp 2-1 \rp / \lp 2+1 \rp = 1/3$,
		so the fraction of incident power accepted is
		$1 - \vert \Gamma \vert ^2 = 1 - 1/9 = 0.889$:
		\eq{ \text{mismatch loss} = -10\ \mylog{0.889} = 0.51\db }
		The antenna under test delivers $0.51\db$ less power than a matched
		antenna of the same gain would, so the comparison method reads
		$0.51\db$ \textbf{low}: $31.8\dbi$ measured against $32.3\dbi$ of true
		gain. Report the measured figure as \emph{realized gain} (which is the
		honest system number), and quote the mismatch-corrected $32.3\dbi$
		separately if the customer wants gain proper.
		\end{solution}

\end{parts}

\newpage

\question \textbf{Near-field scanning and the cut list.} The last two pieces of
the midterm: how you measure when the range cannot exist, and what a complete
characterization actually contains.
\begin{parts}

	\part A near-field scanner must record both amplitude \emph{and} phase at
	every sample point. Explain in one or two sentences why magnitude alone will
	not do.
		\begin{solution}[1.0in]
		The near-field to far-field step is a Fourier transform, and a Fourier
		transform operates on the complex field. Magnitude alone does not
		determine the transform -- infinitely many different phase distributions
		share the same magnitude on the scan plane and produce entirely different
		far-field patterns -- so a scanner needs a phase-coherent receiver and
		phase-stable cabling.
		\end{solution}

	\part Explain, in two or three sentences, why the near-field to far-field
	transform is trustworthy rather than a curve fit. Name the earlier result it
	rests on.
		\begin{solution}[1.2in]
		L6 established that the far field is the Fourier transform of the source
		distribution, with $k\sinp{\theta}$ acting as the spatial frequency
		variable. Fourier transforms are invertible, so the complex tangential
		field sampled on a surface enclosing the antenna determines the field
		everywhere outside that surface -- uniquely, with no fitting or
		assumption involved. The transform is simply the same relationship that
		produced the pattern in the first place, run in the other direction.
		\end{solution}

	\begin{minipage}[t]{0.58\linewidth}
		\part A planar scan at $18\ \ghz$ covers a $1.5\ \m$ by $1.5\ \m$ plane
		at the coarsest spacing the sampling rule allows. Find the required
		sample spacing and the total number of sample points.
		\begin{solution}[1.4in]
		\eq{ \lambda &= \frac{3\ee{8}}{18\ee{9}} = 0.01667\ \m \\
		     \Delta &= \frac{\lambda}{2} = 8.33\ee{-3}\ \m = 8.33\ \text{mm} }
		Along one edge there are $1.5 / 8.33\ee{-3} = 180$ intervals, hence 181
		points:
		\eq{ N = \lp 181 \rp^2 = 32{,}761 \approx 3.3\ee{4}\ \text{points} }
		At even a few points per second this is hours of scanning, which is the
		real price of a near-field range.
		\end{solution}
	\end{minipage}\hfill%
	\begin{minipage}[t]{0.37\linewidth}
		\raggedleft \vspace{12pt}
		$\Delta$ = \ansbox[1.1in]{$8.33\ \text{mm}$}
		\\[10pt]
		$N$ = \ansbox[1.1in]{$32{,}761$}
	\end{minipage}

	\part You are handed a linearly polarized pyramidal horn and told to
	characterize it. List the pattern cuts and polarization configurations that
	make up a minimum complete characterization, state how many sweeps that is,
	and name one thing that set of sweeps still will not tell you.
		\begin{solution}[1.5in]
		Two principal-plane great-circle cuts, each run in two polarization
		configurations:
		\begin{enumerate}\itemsep0pt
			\item E-plane, co-polarized (source aligned with the horn's E-field)
			\item E-plane, cross-polarized (source rotated $90\mydeg$)
			\item H-plane, co-polarized
			\item H-plane, cross-polarized
		\end{enumerate}
		That is \textbf{four sweeps}. What it will not tell you is the pattern in
		the diagonal ($45\mydeg$) planes, which for a rectangular aperture is
		precisely where the highest sidelobes usually sit; a conical cut or a
		full spherical scan is needed for those.
		\end{solution}

\end{parts}

\vspace{1.5em}
\noindent\textbf{Documentation:} \rule[-2pt]{0.6\linewidth}{0.4pt}

\end{questions}
